\begin{explanationbox}

	\textbf{\textcolor{CExplanation}{一句话直觉}}

	椭圆、双曲线、抛物线的标准方程分别是各自定义（距离和、差、等距）的直接代数表达；看到方程形式就能对应定义。

	\medskip
	\textbf{\textcolor{CExplanation}{核心拆解}}
	\begin{description}[style=nextline, leftmargin=2.8em, labelsep=0.5em, itemsep=0.45em]
		\item[\textbf{要点一（定义代数化）：}] 根据定义列出含根号的等式，通过平方消去根号，得到标准方程。
		\item[\textbf{要点二（参数几何意义）：}] 椭圆中 $a$ 是半长轴，$c$ 是半焦距，$b=\sqrt{a^2-c^2}$；双曲线中 $a$ 是半实轴，$c$ 是半焦距，$b=\sqrt{c^2-a^2}$；抛物线中 $p$ 是焦点到准线的距离。
		\item[\textbf{要点三（焦点决定形式）：}] 焦点在 $x$ 轴时，椭圆、双曲线分母对应 $x^2$ 的项在第一个，抛物线方程中 $x$ 是一次项；焦点在 $y$ 轴时形式互换。
	\end{description}

	\medskip
	\textbf{\textcolor{CExplanation}{几何本质（轨迹定义）}}

	椭圆 $\{P\mid |PF_1|+|PF_2|=2a\}$ 的方程是 $\frac{x^2}{a^2}+\frac{y^2}{b^2}=1$；双曲线 $\{P\mid ||PF_1|-|PF_2||=2a\}$ 的方程是 $\frac{x^2}{a^2}-\frac{y^2}{b^2}=1$；抛物线 $\{P\mid |PF|=d(P,l)\}$ 的方程是 $y^2=2px$。方程中的每一项都反映了几何对称性。

	\medskip
	\textbf{\textcolor{CExplanation}{代数意义（对称与系数）}}

	方程只含 $x^2,y^2$ 项（抛物线含一次项），不含交叉项，说明曲线关于坐标轴对称；系数 $a,b,c,p$ 确定了曲线的开口大小、位置和形状。

	\medskip
	\textbf{\textcolor{CExplanation}{考点价值（标准应用）}}
	\begin{itemize}[leftmargin=*, itemsep=0.5em, topsep=0.4em]
		\item \textbf{考法一（由方程求参数）：} 给定标准方程，直接读出 $a^2,b^2$ 或 $p$，进而计算焦点、准线、离心率等。
		\item \textbf{考法二（待定系数求方程）：} 根据已知条件（如焦点位置、焦距、曲线上点坐标），设出标准方程形式，代入计算。
	\end{itemize}

	\medskip
	\textbf{\textcolor{CExplanation}{顿悟点}}

	\textbf{距离公式与平方消元}是解析几何中由几何条件推导轨迹方程的通用工具。圆锥曲线标准方程的推导正是这个方法的典型示范。

	\medskip
	\textbf{\textcolor{CExplanation}{使用场景}}
	\begin{itemize}[leftmargin=*, itemsep=0.5em, topsep=0.4em]
		\item \textbf{场景一（图像识别）：} 给定方程，快速判断曲线类型、焦点位置、开口方向。
		\item \textbf{场景二（轨迹求解）：} 已知动点满足圆锥曲线定义，直接套用标准方程形式求轨迹。
	\end{itemize}
\end{explanationbox}
